\indent ವಿಭಕ್ತ್ಯಂತವಾಗಿರುವುದು ; ಹಾಗೆ ಪೂರ್ವಪದವು ಯಾವ ವಿಭಕ್ತಿಯಲ್ಲಿದೆ ಯೋ ಆ ವಿಭಕ್ತಿಯ ಹೆಸರಿನೊಂದಿಗೆ ಸಮಾಸ ಹೇಳುವುದು ರೂಢಿಯಾಗಿದೆ. ಮೇಲಿನ ಉದಾಹರಣೆಯಲ್ಲಿ “ತಲೆಯಲ್ಲಿ” ಎಂಬ ಪದವು ಸಪ್ತಮೀ ವಿಭಕ್ತಿಯಲ್ಲಿದೆ. ಆದುದರಿಂದ “ತಲೆನೋವು” ಎಂಬ ಶಬ್ದವು ಸಪ್ತಮೀ ತತ್ಪುರುಷ ಸಮಾಸ (ಸ.ತ.ಸ.) ವಾಗಿರುವುದು.\par

\vspace{3pt plus 2pt}\noindent\hspace*{1.2em}\kbvlabel{ಉದಾ}\par\vspace{1pt}

\vspace{5pt plus 4pt minus 1pt}
\noindent{\renewcommand{\arraystretch}{1.24}\setlength{\tabcolsep}{4pt}\rowcolors{1}{kbvrow}{white}
\begin{tabularx}{\textwidth}{X X X X X X}
ಮುಳ್ಳಿನ &  +  & ಹಾಸಿಗೆ &  =  & ಮುಳ್ಳುಹಾಸಿಗೆ & (ಸ.ತ.ಸ.) \\
ತಲೆಯ &  +  & ಕೂದಲು &  =  & ತಲೆಗೂದಲು & (ಸ.ತ.ಸ.) \\
ಮಳೆಯ &  +  & ಕಾಲ &  =  & ಮಳೆಗಾಲ & (ಸ.ತ.ಸ.) \\
ದೇವರ &  +  & ಮಂದಿರ &  =  & ದೇವಮಂದಿರ & (ಷ.ತ.ಸ.) \\
ವಯಸ್ಸಿನಿಂದ &  +  & ವೃದ್ಧ &  =  & ವಯೋವೃದ್ಧ & (ತೃ.ತ.ಸ.) \\
ವ್ಯಾಘ್ರದ &  +  & ದೆಸೆಯಿಂದ &  =  & ವ್ಯಾಘ್ರಭಯ & (ಪಂ.ತ.ಸ.) \\
ನರರಲ್ಲಿ &  +  & ಉತ್ತಮ &  =  & ನರೋತ್ತಮ & (ಸ.ತ.ಸ.) \\
\end{tabularx}}
\vspace{5pt plus 4pt minus 1pt}

\vspace{4pt plus 4pt minus 1pt}\noindent{\kbvsecfont\fontsize{11pt}{14.5pt}\selectfont\color{kbvaccentdark} (೨)\hspace{0.45em}ಕರ್ಮಧಾರಯ ಸಮಾಸ :}\hspace{0.5em}ಈ ವಿಗ್ರಹ ವಾಕ್ಯಗಳನ್ನು ನೋಡೋಣ.\par

\vspace{5pt plus 4pt minus 1pt}
\noindent{\renewcommand{\arraystretch}{1.24}\setlength{\tabcolsep}{4pt}
\begin{tabularx}{\textwidth}{X X X X X X}
ದೊಡ್ಡವನು &  +  & ಅಪ್ಪನು &  =  & ದೊಡ್ಡಪ್ಪನು & (ಪು.ಬ.) \\
ಹೆತ್ತವಳು &  +  & ಅಮ್ಮನು &  =  & ಹೆತ್ತಮ್ಮನು & (ಸ್ತ್ರೀ.ಬ.) \\
ಚಿಕ್ಕವು &  +  & ಮಕ್ಕಳು &  =  & ಚಿಕ್ಕಮಕ್ಕಳು & (ನ.ಬ.) \\
\end{tabularx}}
\vspace{5pt plus 4pt minus 1pt}

\indent ಈ ಮೂರು ಉದಾಹರಣೆಗಳಲ್ಲಿಯೂ ಪೂರ್ವ ಮತ್ತು ಉತ್ತರಪದ ಗಳಲ್ಲಿ ಲಿಂಗ ಭೇದವಾಗಲೀ ವಚನ ಭೇದವಾಗಲೀ ವಿಭಕ್ತಿರೂಪ ಭೇದವಾಗಲೀ ಕಂಡುಬರುವುದಿಲ್ಲ. ಅಪ್ಪನು, ಅಮ್ಮನು, ಮಕ್ಕಳು ಉತ್ತರಪದಗಳಾಗಿದ್ದು\par

\indent ಅವುಗಳ ಅರ್ಥವೇ ಪ್ರಧಾನವಾಗಿರುವುದು. ಅದೇ ಅಲ್ಲದೇ ದೊಡ್ಡವನು, ಹೆತ್ತವಳು, ಚಿಕ್ಕವು ಎಂಬ ಪದಗಳು ಕ್ರಮವಾಗಿ ಅಪ್ಪನು ಅಮ್ಮನು, ಮಕ್ಕಳು ಎಂಬ ಪದಗಳಿಗೆ ವಿಶೇಷಣಗಳಾಗಿವೆ. ಅಂದರೆ, ಮೇಲೆ ಕೊಟ್ಟಿರುವ ಪ್ರತಿಯೊಂದು ಉದಾಹರಣೆಯಲ್ಲಿ ಪೂರ್ವಪದವು ವಿಶೇಷಣ ಮತ್ತು ಉತ್ತರ\par

\indent ಪದವು ವಿಶೇಷ್ಯ ಆಗಿರುವುದು. ಇವು ಕರ್ಮಧಾರಯ ಸಮಾಸದ ಲಕ್ಷಣಗಳು.\par
